% !TeX program = lualatex
% Compilar desde esta carpeta: lualatex -output-directory=build -interaction=nonstopmode -halt-on-error Informe_datos_proyecto.tex
\documentclass[11pt,a4paper,oneside]{report}
\usepackage{fontspec}
\setmainfont{Arial}
\usepackage[spanish,es-nodecimaldot,es-noshorthands]{babel}
\usepackage{geometry}
\geometry{a4paper,left=3cm,top=3cm,right=2.5cm,bottom=2.5cm,headheight=14pt,headsep=0.6cm}
\usepackage{setspace}
\onehalfspacing
\usepackage{graphicx}
\usepackage{booktabs}
\usepackage{longtable}
\usepackage{tabularx}
\usepackage{array}
\usepackage{pdflscape}
\usepackage{ragged2e}
\usepackage{xcolor}
\usepackage{colortbl}
\usepackage{fancyhdr}
\usepackage{titlesec}
\usepackage[hidelinks]{hyperref}
\usepackage{bookmark}
\usepackage{enumitem}
\usepackage{microtype}
\usepackage{tikz}
\usetikzlibrary{arrows.meta,positioning}
\usepackage{pifont}
\definecolor{unired}{HTML}{8B1A1A}
\definecolor{darkblue}{HTML}{17365D}
\definecolor{lightblue}{HTML}{DCE6F1}
\definecolor{baseline}{HTML}{5A7D71}
\definecolor{active}{HTML}{E9B44C}
\definecolor{planned}{HTML}{DCE6F1}
\setlength{\parindent}{0pt}
\setlength{\parskip}{9pt}
\setlength{\emergencystretch}{3em}
\setlist{leftmargin=1.1cm,itemsep=0.35\baselineskip,topsep=0.35\baselineskip}
\titleformat{\chapter}[hang]{\normalfont\bfseries\fontsize{11}{13.2}\selectfont}{\MakeUppercase{\chaptername}\ \thechapter.}{0.75em}{\MakeUppercase}
\titleformat{name=\chapter,numberless}[block]{\normalfont\bfseries\fontsize{11}{13.2}\selectfont}{}{0pt}{\MakeUppercase}
\titleformat{\section}{\normalfont\bfseries\fontsize{11}{13.2}\selectfont}{\thesection}{0.75em}{\MakeUppercase}
\titleformat{\subsection}{\normalfont\bfseries\fontsize{11}{13.2}\selectfont}{\thesubsection}{0.75em}{\MakeUppercase}
\titlespacing*{\chapter}{0pt}{0pt}{0.75\baselineskip}
\fancypagestyle{plain}{\fancyhf{}\fancyhead[R]{\thepage}\renewcommand{\headrulewidth}{0pt}}
\pagestyle{plain}
\newcommand{\autores}{LUIS ENRIQUE KOC G\'ONGORA\\HERBERT ANTONIO MEL\'ENDEZ GARC\'IA}
\newcommand{\asesor}{Samuel Oporto}
\newcommand{\caratula}[1]{%
  \hypersetup{pageanchor=false}
  \begin{titlepage}
  \thispagestyle{empty}
  \begin{center}
  {\bfseries\fontsize{19}{23}\selectfont UNIVERSIDAD NACIONAL DE INGENIER\'IA\par}
  \vspace{0.2cm}
  {\bfseries\fontsize{13}{16}\selectfont UNIDAD DE POSGRADO DE LA FACULTAD DE INGENIER\'IA INDUSTRIAL Y DE SISTEMAS\par}
  {\bfseries\fontsize{13}{16}\selectfont ESCUELA DE POSGRADO\par}
  \vspace{0.25cm}
  \includegraphics[width=4.6cm]{../Plan/imagenes/logo_UNI.png}\par
  \vspace{0.2cm}
  {\bfseries\fontsize{14}{17}\selectfont #1\par}
  \vspace{0.35cm}
  {\bfseries\fontsize{12}{15}\selectfont DISE\~NO Y EVALUACI\'ON DE UNA RED NEURONAL INFORMADA POR F\'ISICA PARA ESTIMAR LA TEMPERATURA Y LA AMPACIDAD DE CABLES EL\'ECTRICOS ENTERRADOS EN ENTORNOS HETEROG\'ENEOS\par}
  \vfill
  {\fontsize{11}{14}\selectfont PARA OBTENER EL GRADO ACAD\'EMICO DE\\MAESTRO EN CIENCIAS CON MENCI\'ON EN INTELIGENCIA ARTIFICIAL\par}
  \vspace{0.35cm}
  {\bfseries\fontsize{11}{14}\selectfont PRESENTADO POR:\par}
  {\bfseries\fontsize{11}{14}\selectfont \autores\par}
  \vspace{0.3cm}
  {\fontsize{11}{14}\selectfont ASESOR: \asesor\par}
  \vfill
  {\fontsize{11}{14}\selectfont Lima, Per\'u\\2026\par}
  \end{center}
  \end{titlepage}
  \hypersetup{pageanchor=true}
  \pagenumbering{arabic}
  \tableofcontents
  \clearpage
}
\begin{document}
\caratula{INFORME DE DATOS, FUENTES Y CASOS DE ESTUDIO}

\chapter*{Resumen ejecutivo}
\addcontentsline{toc}{chapter}{Resumen ejecutivo}
El conjunto documental de partida contiene 16 instancias de estudio (nueve f\'isicas o de campo y siete de verificaci\'on) y el cat\'alogo activo de PINN\_v2 contiene 21 definiciones aceptadas por su cargador. Este informe conserva ese inventario, incorpora los dos casos de cable de CIGRE TB 880 y TB 963 adem\'as del anillo anal\'itico, y ampl\'ia la evidencia peruana con proyectos COES y segmentos subterr\'aneos que aparecen en el Anuario Estad\'istico de Electricidad 2024 del MINEM.

La variable de entrada del proyecto es la descripci\'on geom\'etrica de la instalaci\'on enterrada y de su entorno: capas del cable, arreglo, ductos o zanja, profundidad, estratos, relleno, propiedades t\'ermicas del suelo y condiciones ambientales/fronterizas. Las fuentes peruanas consultadas confirman instalaciones o propuestas a distintas tensiones y aportan rutas, longitudes, secciones de conductor y, en algunos registros, capacidad. No informan por s\'i solas todas las dimensiones de cable, configuraci\'on de zanja, profundidad y propiedades t\'ermicas del terreno; por tanto, son evidencia de contexto y candidatos para solicitar datos, no casos PINN listos para ejecutar.

\section*{Objeto, alcance y unidad de registro}
La unidad de datos es un caso o segmento documentado, con identificador estable, fuente, ubicaci\'on exacta, valores y unidades tal como se publicaron, transformaciones aplicadas, estado de procedencia y faltantes. Se distinguen: (a) instalaci\'on reportada/operativa; (b) proyecto de transmisi\'on planificado; (c) caso experimental, normativo o num\'erico publicado; y (d) benchmark anal\'itico o manufacturado. Que un registro pertenezca a (a) o (b) no significa que sea model-ready.

\section*{Criterios y estrategia de b\'usqueda}
Se revisaron el informe de datos previo y su CSV; las fuentes citadas en el plan y el repositorio Zotero local; los PDF locales de CIGRE TB 880 y TB 963; el diagn\'ostico COES 2023 del SEIN 2025--2034; el texto extra\'ido del Anuario MINEM 2024; y los JSON activos de PINN\_v2. Para COES, el cribado abarc\'o referencias a ``subterr\'aneo'', ``l\'inea subterr\'anea'', ``cable'', enlaces y proyectos de transmisi\'on sin filtro de tensi\'on. Se registr\'o tensi\'on cuando estuvo disponible, sin exigir 33 o 60 kV.

Se incluyeron casos cuando la fuente permiti\'o identificar una instalaci\'on/caso, una entrada de verificaci\'on o una variable del entorno pertinente. Se excluyeron de la categor\'ia de model-ready los casos sin geometr\'ia o propiedades t\'ermicas necesarias. No se completaron valores ausentes con promedios ni se infirieron profundidad, formaci\'on, capas o suelo a partir de tensi\'on, secci\'on del conductor o ampacidad.

\chapter*{Inventario previo: 16 instancias}
\addcontentsline{toc}{chapter}{Inventario previo: 16 instancias}
La Tabla~\ref{tab:inventario16} reproduce las 16 claves del \texttt{04\_indice\_instancias.csv} de la revisi\'on previa y conserva su clasificaci\'on y uso propuesto. Para las variables y resultados num\'ericos detallados de Kim, Khumalo, Al-Dulaimi, Atoccsa, Oc\~lo\'n, Aras, Quan, Oladunjoye, M\"obius, CIGRE, IEC, Pan, Hahn y la suite MMS, se conserva como documento de detalle el informe previo de 23 p\'aginas; el libro Excel acompa\~nante facilita la b\'usqueda transversal de las 16 entradas.

\begin{landscape}
\small
\setlength{\tabcolsep}{3pt}
\begin{longtable}{p{2.0cm}p{2.8cm}p{3.8cm}p{4.8cm}}
\caption{Inventario previo de instancias y su utilidad en el proyecto. Fuente: \texttt{04\_indice\_instancias.csv}.}\label{tab:inventario16}\\
\toprule
\textbf{Clave} & \textbf{Fuente / tipo} & \textbf{Instancia documentada} & \textbf{Utilidad propuesta} \\
\midrule
\endfirsthead
\toprule
\textbf{Clave} & \textbf{Fuente / tipo} & \textbf{Instancia documentada} & \textbf{Utilidad propuesta} \\
\midrule
\endhead
kim2025 & Kim et al. (2025); f\'isico & UPCS 154 kV, seis cables, PAC/arena y tres estratos & Campo t\'ermico, ampacidad y bedding \\
khumalo2025 & Khumalo et al. (2025); f\'isico & Cable MV XLPE, secado de suelo y ampacidad & Sensibilidad a humedad y resistividad \\
aldulaimi2024 & Al-Dulaimi et al. (2024); f\'isico & 132 kV, suelo multicapa, SCMB/FTB y FEM--BPNN & Geometr\'ia, materiales y heterogeneidad \\
atoccsa2024 & Atoccsa et al. (2024); f\'isico & 220 kV, backfill t\'ermico y optimizaci\'on PSO & Ampacidad y dise\~no de relleno \\
oclon2015 & Oc\~lo\'n et al. (2015); f\'isico & 400 kV, FTB/SBM/HDPE y PSO & Interfaces y relleno localizado \\
aras2005 & Aras et al. (2005); f\'isico & 154 kV, contraste IEC/FEM/ensayo & Verificaci\'on en suelo homog\'eneo \\
quan2019 & Quan et al. (2019); f\'isico & Tres cables; formaci\'on plana/tr\'ebol; suelo/FTB & Geometr\'ia, profundidad y $k(T)$ \\
oladunjoye2012 & Oladunjoye et al. (2012); campo & Diez calicatas con propiedades del suelo & Parametrizaci\'on de suelo observada \\
mobius2025 & M\"obius et al. (2025); f\'isico & Cables MV en Hamburgo y 32 a\~nos de datos & Escenarios estacionales de suelo \\
cigre2022\_case0 & CIGRE TB 880; normativo & Cable 132 kV y 14 variantes IEC & Regresi\'on corriente--p\'erdidas--temperatura \\
cigre2025\_case1 & CIGRE TB 963; FEM 2D & Tres cables 132 kV en ductos, nueve variantes & Campo, dominio, interfaces y ampacidad \\
iec60853\_annexA & IEC 60853-3; transitorio & Cable tripolar 132 kV; ciclo y secado parcial & Extensi\'on temporal fuera del n\'ucleo estacionario \\
pan2025\_plate & Pan et al. (2025); PINN 2D & Placa rectangular de Laplace y FEM & Muestreo e inversi\'on de frontera \\
cigre2025\_annulus & CIGRE TB 963; anal\'itico & Anillo con flujo interior y temperatura exterior & Error de campo contra soluci\'on exacta \\
hahn2012\_rect & Hahn y \"Ozi\c{s}ik; anal\'itico & Rect\'angulo 2D con frontera convectiva & Validaci\'on de Robin \\
mms2d\_suite & Elaboraci\'on propia; manufacturado & Conductividad variable e interfaz multimaterial & Pruebas de PDE, $k(x,y)$ e interfaces \\
\bottomrule
\end{longtable}
\end{landscape}

\chapter*{Fuentes CIGRE y referencias de verificaci\'on}
\addcontentsline{toc}{chapter}{Fuentes CIGRE y referencias de verificaci\'on}
\textbf{CIGRE TB 880, caso introductorio 0 (2022, pp. 61--125 del informe previo).} Cable 132 kV, tres conductores unipolares Cu de 630 mm\textsuperscript{2} en formaci\'on trebolada, directamente enterrados con centro a 1,0 m; suelo a 20 \textdegree C y resistividad t\'ermica 1,0 K m/W. El caso base reporta ampacidades convergidas de 821,78 A, 886,18 A y 803,16 A para tratamientos de p\'erdidas/pantallas diferenciados. La familia incluye instalaci\'on directa, ductos HDPE y PVC en concreto, aire con radiaci\'on y canal. Es valioso para la cadena normativa IEC y variaciones de instalaci\'on; no entrega por s\'i solo un campo FEM 2D.

\textbf{CIGRE TB 963, caso 1 (2025, pp. 103--121 del informe previo).} Tres cables de 132 kV y 630 mm\textsuperscript{2} en ductos HDPE horizontales dentro de backfill, rodeados por suelo y asfalto. El documento reporta nueve variantes FEM 2D; incluye dominio, capas del cable, ductos, temperatura, resistividades, ampacidades y sensibilidad. Es la referencia publicada m\'as directa para comparar campos, interfaces y resultados t\'ermicos del dominio 2D.

\textbf{CIGRE TB 963, ejemplo anular (2025, pp. 52--55).} Caso anal\'itico de conducci\'on con radio interior 0,06 m, exterior 20 m, resistividad 0,8 K m/W, borde exterior a 20 \textdegree C y flujo interno equivalente a 40 W/m. Se conserva como verificaci\'on exacta; no sustituye los dos casos de cable anteriores.

\chapter*{Instalaciones y proyectos peruanos}
\addcontentsline{toc}{chapter}{Instalaciones y proyectos peruanos}
\textbf{COES, diagn\'ostico 2025--2034.} En la tabla de proyectos de transmisi\'on que acompa\~na el informe COES/DP-01-2023 (p. 69 impresa) figura el ``Enlace Subterr\'aneo 220 kV Chavarr\'ia--San Mart\'in'', ligado a una nueva SE San Mart\'in y acometidas subterr\'aneas. Es evidencia de proyecto planificado, no una instalaci\'on ya ejecutada ni una ficha de dise\~no t\'ermico. En la discusi\'on operativa del mismo diagn\'ostico se describen enlaces subterr\'aneos de 220 kV Manchay--Pachac\'utec y Pachac\'utec--Progreso. El PDF separado de la propuesta definitiva COES 2023--2032 se cita en la bibliograf\'ia del plan, pero no se obtuvo desde su URL durante esta preparaci\'on (la descarga devolvi\'o un archivo vac\'io); por ello no se atribuyen a ese PDF hallazgos no inspeccionados.

\textbf{MINEM, Anuario Estad\'istico de Electricidad 2024.} El texto extra\'ido contiene 139 filas con la palabra ``Subterr\'aneo''; el recuento es de menciones de texto, no de 139 activos \emph{\'unicos}. El Excel preserva las 139 filas crudas con localizador, y presenta adem\'as una selecci\'on normalizada por c\'odigo. Entre los ejemplos verificables est\'an L-2005/L-2006 Chavarr\'ia--Barsi (1,00 km, 220 kV, 335 MVA, Cu 1200 mm\textsuperscript{2}); L-2114 y L-2115 Mirador--Malvinas (cada uno 4,08 km, 220 kV, Cu 1200 mm\textsuperscript{2}); L-613 Santa Rosa Antigua--Tacna (2,95 km, 60 kV, 62 MVA, Cu 400 mm\textsuperscript{2}); L-621/L-622 Chavarr\'ia--Mirones (1,89 km, 60 kV, 56 MVA, Cu 500 mm\textsuperscript{2}); L-6349 Huarangal--Comas (7,84 km, 60 kV, capacidad no consignada, Cu 1200 mm\textsuperscript{2}); y L-6556 Chavarr\'ia--UNI (5,00 km, 60 kV, 67 MVA, Cu 500 mm\textsuperscript{2}). Son registros de inventario y contexto peruano. No proporcionan profundidad, arreglo, secci\'on de zanja, geometr\'ia de ductos ni propiedades del suelo, de modo que no son por s\'i mismos casos de simulaci\'on.

\textbf{Arias Velasquez et al. (2023).} El art\'iculo estudia el dise\~no de circuitos subterr\'aneos de media tensi\'on y la influencia t\'ermica; es una fuente cercana para criterios de dise\~no, pero no debe confundirse con el inventario de una l\'inea espec\'ifica del SEIN. Atoccsa et al. (2024) es otra contribuci\'on de autores peruanos sobre cable de 220 kV y backfill optimizado; se conserva como caso bibliogr\'afico/de dise\~no, no como evidencia de que esa geometr\'ia corresponda a una instalaci\'on peruana en operaci\'on.

\chapter*{Modelo de datos, procedencia y preprocesamiento}
\addcontentsline{toc}{chapter}{Modelo de datos, procedencia y preprocesamiento}
\small
\setlength{\tabcolsep}{3pt}
\begin{longtable}{p{2.8cm}p{5.1cm}p{5.8cm}}
\toprule
\textbf{Grupo de entrada} & \textbf{Campos que se deben conservar} & \textbf{Uso y controles} \\
\midrule
Geometr\'ia del cable & Tipo de conductor, radios/di\'ametros, espesores, secuencia de capas & Convertir unidades a SI sin perder valor original; comprobar radios crecientes \\
Geometr\'ia de instalaci\'on & Coordenadas, arreglo plano/tr\'ebol, n\'umero de cables, profundidad, ductos, zanja, backfill y dominio & No inferir formaci\'on ni profundidad si no se publican \\
Suelo y entorno & Capas/estratos, clasificaci\'on, $k$, resistividad, humedad, densidad, temperatura, frontera superficial & Distinguir propiedades publicadas, extra\'idas de tabla, escenarios inferidos y supuestos \\
Operaci\'on y fuente & Tensi\'on, corriente, resistencia, p\'erdidas, factor de carga & Tensi\'on y capacidad de inventario no equivalen a fuente volum\'etrica del caso PINN \\
Referencia y resultados & Campo FEM/anal\'itico, $T_{\max}$, ampacidad, malla y criterios de convergencia & Separar resultado publicado de predicci\'on PINN; declarar error de interpolaci\'on \\
\bottomrule
\end{longtable}

El preprocesamiento previsto conserva el valor transcrito, unidad original y valor normalizado; documenta p\'agina/tabla y nombre del archivo; asigna un estado \texttt{publicado}, \texttt{extraido}, \texttt{inferido}, \texttt{supuesto} o \texttt{faltante}; y registra cada transformaci\'on, adaptaci\'on e inconsistencia. Los estados heredados del cat\'alogo tambi\'en se reportan: 20 de las 21 definiciones activas llevan al menos un grupo marcado como \texttt{supuesto}; aunque todas pasan actualmente el cargador, eso no equivale a una auditor\'ia de procedencia por campo.

El CSV previo presenta diferencias tipogr\'aficas en algunos DOI respecto de la bibliograf\'ia, por lo que el DOI y el t\'itulo se cotejan con \texttt{Plan/referencias\_plan\_tesis.bib}; por ejemplo, Atoccsa et al. (2024) aparece como \texttt{10.3390/en17051023} en la bibliograf\'ia de control. Los duplicados del MINEM se mantienen como registros separados cuando el c\'odigo de circuito o longitud distingue segmentos, y se vinculan mediante una clave normalizada en Excel; no se suman como una sola instalaci\'on sin soporte documental.

\section*{L\'imites de uso}
\begin{itemize}
  \item El reporte previo proporciona valores detallados de literatura y benchmarks; los inventarios COES/MINEM aportan prueba de contexto peruano, pero no las propiedades t\'ermicas y geom\'etricas necesarias para una PINN f\'isica completa.
  \item El COES 2022 definitivo se localiz\'o por cita bibliogr\'afica; no se pudo inspeccionar el PDF separado en esta ejecuci\'on. El proyecto Chavarr\'ia--San Mart\'in est\'a apoyado por la menci\'on del diagn\'ostico COES local, no por una lectura directa del archivo de 2022.
  \item Los casos JSON incompletos se incluir\'an en el paquete como candidatos no ejecutables. La versi\'on actual del cargador rechaza datos requeridos marcados como faltantes; no se altera el esquema para simularlos.
\end{itemize}

\chapter*{Referencias principales}
\addcontentsline{toc}{chapter}{Referencias principales}
Al-Dulaimi, A. A., et al. (2024). Adaptive FEM-BPNN model for predicting underground cable temperature considering varied soil composition. \textit{Engineering Science and Technology, an International Journal, 51}, 101658. https://doi.org/10.1016/j.jestch.2024.101658\\
Arias Velasquez, R. M., Gonzales-Condori, E. G., \& Velarde Loayza, P. E. (2023). Enhancing medium voltage underground circuit design: Assessing limitations, thermal influence, and accurate modelling. \textit{Results in Engineering, 20}, 101552. https://doi.org/10.1016/j.rineng.2023.101552\\
Atoccsa, B. A., et al. (2024). Optimization of ampacity in high-voltage underground cables with thermal backfill using dynamic PSO and adaptive strategies. \textit{Energies, 17}(5), 1023. https://doi.org/10.3390/en17051023\\
CIGRE Working Group B1.56. (2022). \textit{Power cable rating examples for calculation tool verification} (Technical Brochure 880).\\
CIGRE Working Group B1.87. (2025). \textit{Finite element analysis for cable rating calculations} (Technical Brochure 963).\\
COES. (2023). \textit{Informe de diagn\'ostico de las condiciones operativas del SEIN, periodo 2025--2034} (Informe COES/DP-01-2023).\\
Ministerio de Energ\'ia y Minas. (2025). \textit{Anuario Estad\'istico de Electricidad 2024}. Tabla de l\'ineas de transmisi\'on; fuente local en \texttt{Plan/datos\_fuente/descargas/minem\_anuario\_2024.pdf} y texto extra\'ido.
\end{document}